\chapter*{Mở đầu}
\setcounter{page}{1}
\pagenumbering{arabic}

Trong bối cảnh thiết bị di động trở thành công cụ đồng hành thường xuyên trong đời sống hằng ngày, các ứng dụng sức khỏe và thể hình ngày càng đóng vai trò quan trọng trong việc hỗ trợ người dùng xây dựng thói quen sinh hoạt tích cực. Tuy nhiên, đối với nhóm người tập gym, việc duy trì một quá trình luyện tập hiệu quả không chỉ phụ thuộc vào số lần đến phòng tập mà còn đòi hỏi phải có kế hoạch tập phù hợp, hiểu đúng kỹ thuật bài tập, theo dõi sự thay đổi của cơ thể và kiểm soát chế độ dinh dưỡng theo mục tiêu tăng cơ hoặc giảm mỡ.

Trên thực tế, nhiều ứng dụng hiện nay mới chỉ hỗ trợ một phần rời rạc của bài toán, chẳng hạn như thư viện bài tập, bộ đếm calo, hoặc lịch nhắc vận động. Điều đó khiến người dùng phải thao tác trên nhiều hệ thống khác nhau để hoàn thành một quy trình thống nhất từ xác định mục tiêu, lựa chọn bài tập, duy trì lịch tập đến theo dõi kết quả dài hạn. Đặc biệt, người mới bắt đầu thường gặp khó khăn khi phải tự xây dựng lộ trình tập luyện do thiếu kiến thức chuyên môn và thiếu công cụ đủ nhất quán để đồng hành.

Từ nhu cầu đó, đề tài hướng tới việc xây dựng \textbf{Fitty}, một ứng dụng di động hỗ trợ tập gym và theo dõi sức khỏe cá nhân hóa trên nền tảng Android. Ý tưởng cốt lõi của hệ thống là xây dựng một trợ lý số dành cho người tập luyện, giúp họ bắt đầu từ hồ sơ cá nhân ban đầu, nhận kế hoạch tập khởi đầu phù hợp, truy cập thư viện bài tập có hướng dẫn trực quan, thực hiện buổi tập theo lịch, ghi nhận bữa ăn, theo dõi các chỉ số như cân nặng, BMI và tỷ lệ mỡ cơ thể ước lượng, đồng thời duy trì động lực qua cơ chế nhắc việc và thông báo.

Để hiện thực mục tiêu đó, báo cáo tập trung vào việc phân tích yêu cầu, thiết kế kiến trúc, tổ chức mô hình dữ liệu và triển khai nguyên mẫu chức năng cho các phân hệ chính của hệ thống. Các chương tiếp theo lần lượt trình bày bối cảnh và bài toán, phương án phân tích và thiết kế, lựa chọn công nghệ, quá trình cài đặt và đánh giá kết quả triển khai của Fitty như một ứng dụng hỗ trợ gymer theo hướng cá nhân hóa.

\cleardoublepage
